\documentclass[twocolumn]{aastex631}
\shorttitle{Disentangling aperture and calibration systematics in the MZR}
\shortauthors{NebulaMind Lab}
\begin{document}
\title{Disentangling Aperture and Calibration Systematics in the Gas-Phase Mass--Metallicity Relation: A Practitioner's Framework (Review)}
\author{NebulaMind Lab (autonomous run)}
\affiliation{Descriptive methods review, machine-generated --- not a validated result. Contribution is the synthesis/framework, not a new measurement.}

\begin{abstract}
Comparisons of the gas-phase mass--metallicity relation (MZR) across surveys, simulations, and cosmic
epochs are limited less by statistics than by two systematics that are frequently conflated: the
\emph{calibration-scale} offset between abundance diagnostics, and the \emph{aperture} bias of
fixed-size spectroscopic apertures. We review the recent literature and lay out a practitioner's
framework that separates the two. The calibration-scale offset between theoretical photoionization and
empirical direct-$T_e$ diagnostics reaches up to $\sim0.7$~dex in absolute normalization; aperture bias
introduces errors exceeding $\sim0.15$~dex whenever the fiber covering fraction falls below $\sim20\%$,
and can masquerade as physical structure in the MZR (including its high-mass turnover). Diffuse ionized
gas (DIG), contributing 30--60\% of H$\alpha$, further biases low-ionization line ratios and can flatten
or invert inferred gradients. We recommend: (i) reporting all abundances on a single declared scale with
published conversions; (ii) enforcing covering-fraction floors or aperture corrections calibrated against
integral-field data; and (iii) using resolved (IFS) measurements as the independent ground truth for
integrated relations. This is a synthesis and methods contribution, not a new measurement.
\end{abstract}

\section{Introduction and Scope}
The MZR is a foundational scaling relation used to constrain feedback and chemical-evolution models. Its
interpretation, however, is entangled with observational systematics that do not cancel across samples.
This review addresses two dominant, separable effects --- calibration scale and aperture --- and the role
of diffuse ionized gas, and proposes an explicit framework to keep them distinct in cross-sample work.

\section{Calibration-Scale Systematics}
Direct-$T_e$ abundances are the gold standard but are biased toward metal-poor, high-excitation systems.
Strong-line diagnostics extend coverage but split into theoretical (photoionization) and empirical
($T_e$-anchored) families whose absolute normalizations differ by up to $\sim0.7$~dex --- a factor of five
in metal content. $R_{23}$ is double-valued; $N2$/$O3N2$ implicitly assume a local N/O--O/H scaling that
fails in the N/O-enhanced high-$z$ regime. Mixing diagnostics across redshift can mimic or mask true MZR
evolution. The mitigation is a single declared scale with published inter-calibration conversions.

\section{Aperture Bias}
A fixed angular aperture (e.g.\ the 3\arcsec\ SDSS fiber) samples a redshift-dependent physical radius.
Below a $\sim20\%$ covering fraction, nuclear-to-global metallicity differences exceed $\sim0.15$~dex;
because disks have negative gradients, a small central aperture over-estimates the global abundance and
imprints spurious size--metallicity trends. Massive, extended galaxies remain under-sampled even above the
canonical $z>0.04$ cut, so part of the MZR high-mass turnover is aperture-induced rather than physical.

\section{Diffuse Ionized Gas}
DIG supplies 30--60\% of H$\alpha$ with a harder ionizing field and enhanced low-ionization lines. In
integrated or coarsely-resolved spectra, line-ratio mixing inflates $N2$/$O3N2$ abundances and flattens
--- occasionally inverts --- gradients, which can be misread as pristine-gas infall. Masking
low-$\Sigma_{{\rm H}\alpha}$ regions or using DIG-robust indices ($N2O2$) mitigates this.

\section{A Framework and Recommendations}
(1) Declare one abundance scale and convert all inputs to it. (2) Apply covering-fraction floors or
IFS-calibrated aperture corrections; propagate the residual as a systematic. (3) Treat resolved (IFS) MZR
measurements as the independent ground truth for integrated relations. Applying this separation makes the
$\Delta({\rm O/H})_{\rm aperture}$ term explicitly distinct from the calibration-scale offset --- a
prerequisite for any credible MZR-evolution or simulation-vs-observation claim.

\begin{thebibliography}{}
\bibitem[Andrews \& Martini(2013)]{am13} Andrews, B.~H., \& Martini, P. 2013, ApJ, 765, 140
\bibitem[Kewley et al.(2008)]{kewley08} Kewley, L.~J., \& Ellison, S.~L. 2008, ApJ, 681, 1183
\bibitem[Poetrodjojo et al.(2019)]{dig} Poetrodjojo, H., et al. 2019, MNRAS, 487, 79
\bibitem[S{\'a}nchez et al.(2021)]{ifs} S{\'a}nchez, S.~F. 2021, ARA\&A, 59, 191
\end{thebibliography}
\end{document}
